\subsection{不可约表示的最高权}

对任意表示 \(\tau : \mathrm{G} \to \mathrm{GL}(\mathrm{V})\)，关联同态 \(\mathrm{L}(\tau)_{(\mathbf{C})}\)，它从 \(\mathbf{C}\)-李代数 \(\mathfrak{g}_{\mathbf{C}}\) 到 \(\operatorname{End}(\mathrm{V})\)，并延拓了线性表示 \(\mathrm{L}(\tau)\)（它是 \(\mathfrak{g}\) 在 V 所对应的实向量空间上的表示）（第 III 章，第 3 节，第 11 号）。根据第 4 节第 3 号的命题 7，映射 \(\delta\) 从 \(\mathrm{X}(\mathrm{T})\) 到 \(\operatorname{Hom}_{\mathbf{C}}(\mathfrak{t}_{\mathbf{C}}, \mathbf{C}) = \mathfrak{t}_{\mathbf{C}}^{*}\) 诱导出一个双射：它将 \(\tau\) 关于 T 的权集合双射到 \(\mathrm{L}(\tau)_{(\mathbf{C})}\) 关于 \(\mathfrak{t}_{\mathbf{C}}\) 的权集合，其中后者是 \(\mathfrak{g}_{\mathbf{C}}\) 的嘉当子代数。

引理 2. 设 \(\varphi\) 是复李代数 \(\mathfrak{g}_{\mathbf{C}}\) 在有限维复向量空间 V 上的线性表示。存在 G 在 V 上的表示 \(\tau\)，使得 \(\mathrm{L}(\tau)_{(\mathbf{C})} = \varphi\)，当且仅当 \(\varphi\) 是半单的且 \(\mathbf{t}_{\mathbf{C}}\) 在 V 上的权属于 \(\delta (\mathrm{X(T)})\)。

若存在 G 的表示 \(\tau\) 使得 \(\mathrm{L}(\tau)_{(\mathbf{C})} = \varphi\)，则 \(\varphi\) 是半单的，因为 G 连通且 \(\tau\) 半单（第 III 章，第 6 节，第 5 号，命题 13 的推论 2），并且 \(\mathfrak{t}_{\mathbf{C}}\) 在 V 上的权属于 \(\delta\) 的像。因此该条件是必要的；下面证明其充分性。若 \(\varphi\) 半单，则 V 是各个 \(V_{\mu}(t_{\mathbf{C}})\) 的直和，其中 \(\mu\) 属于 \(\mathfrak{t}_{\mathbf{C}}\) 在 V 上的权集合（第 VII 章，第 2 节，第 4 号，定理 2 的推论 3）；若所有权都属于 \(\delta\) 的像，则存在 T 在 V 上的表示 \(\tau_{\mathrm{T}}\)，使得 \(\mathrm{L}(\tau_{\mathrm{T}})_{(\mathbf{C})} = \varphi |t_{\mathbf{C}}\)：事实上，只需置 \(\tau_{\mathrm{T}}(t)v = t^{\lambda}v\)，其中 \(t \in \mathrm{T}\) 且 \(v \in V_{\delta(\lambda)}(t_{\mathbf{C}})\)。引理现由第 2 节第 6 号命题 8 得到。

定理 1. a) 设 \(\tau : \mathbf{G} \to \mathbf{GL}(\mathbf{V})\) 是 \(\mathbf{G}\) 的不可约表示。则 \(\tau\) 的权集合（相对于 \(\mathbf{T}\)）有一个最大元 \(\lambda\)，它是支配的，并且空间 \(\mathrm{V}_{\lambda}(\mathrm{T})\) 的维数为 1。

\begin{enumerate}
\def\labelenumi{\alph{enumi})}
\setcounter{enumi}{1}
\item
  G 的两个不可约表示等价，当且仅当它们的最高权相等。
\item
  对于 X(T) 的每个支配元素 \(\lambda\)，存在一个最高权为 \(\lambda\) 的 G 的不可约表示。
\end{enumerate}

由引理 2，G 的不可约表示的等价类与 g 的权属于 \(\delta(\mathrm{X}(\mathrm{T}))\) 的有限维不可约表示的等价类之间存在双射。

将 \(\mathcal{G}\mathfrak{g}_{\mathbf{C}}\) 记为中心，将 \(\mathcal{D}\mathfrak{g}_{\mathbf{C}}\) 记为 \(\mathfrak{g}_{\mathbf{C}}\) 的导出李代数，从而 \(\mathfrak{g}_{\mathbf{C}} = \mathcal{G}\mathfrak{g}_{\mathbf{C}} \oplus \mathcal{D}\mathfrak{g}_{\mathbf{C}}\)。对于每个双线性型 \(\mu\)，其定义域为 \(\mathfrak{t}_{\mathbf{C}} \cap \mathcal{D}\mathfrak{g}_{\mathbf{C}}\)，将第 VIII 章第 6 节第 3 号引入的 \(\mathrm{E}(\mu)\) 记为 \(\mathcal{D}\mathfrak{g}_{\mathbf{C}}\)-单模；对于每个线性形式 \(\nu\)，在 \(\mathcal{G}\mathfrak{g}_{\mathbf{C}}\) 上取与之关联的模 \(\mathbf{C}(\nu)\)，它是 \(\mathcal{G}\mathfrak{g}_{\mathbf{C}}\) 上的一维 \(\mathbf{C}\)-模。则 \(\mathfrak{g}_{\mathbf{C}}\)-模 \(\mathbf{C}(\nu) \otimes \mathrm{E}(\mu)\) 是单模；并且根据第 VIII 章第 7 节第 2 号定理 1 的推论 2 以及《代数》第 VIII 章第 11 节第 1 号定理 1，每个有限维 \(\mathfrak{g}_{\mathbf{C}}\)-单模都同构于这些模之一 \(\mathbf{C}(\nu) \otimes \mathrm{E}(\mu)\)；此外（同上引文），\(\mathbf{C}(\nu) \otimes \mathrm{E}(\mu)\) 是有限维的，当且仅当 \(\mu(H_{\alpha})\) 对每个单根 \(\alpha\) 为正整数。若将 \(\nu + \mu\) 记为 \(\mathfrak{t}_{\mathbf{C}}\) 上的线性形式，它诱导 \(\nu\) 于 \(\mathcal{G}\mathfrak{g}_{\mathbf{C}}\)，并诱导 \(\mu\) 于 \(\mathfrak{t}_{\mathbf{C}} \cap \mathcal{D}\mathfrak{g}_{\mathbf{C}}\)，则 \((\nu + \mu)(H_{\alpha}) = \mu(H_{\alpha})\)；此外，\(\mathbf{C}(\nu) \otimes \mathrm{E}(\mu)\) 的权是 \(\nu + \lambda\)，其中 \(\lambda\) 是 \(\mathrm{E}(\mu)\) 的任一权，因而具有 \(\nu + \mu - \theta\) 的形式，其中 \(\theta \in \delta(X_{+})\)（第 VIII 章，第 6 节，第 2 号，引理 2）。

我们由此得出，\(\mathbf{C}(\nu)\otimes\mathrm{E}(\mu)\) 是有限维的，当且仅当 \((\nu+\mu)(H_{\alpha})\) 对每个单根 \(\alpha\) 为正整数；其权属于 \(\delta(\mathrm{X}(\mathrm{T}))\)，当且仅当 \(\nu+\mu\) 属于 \(\delta(\mathrm{X}(\mathrm{T}))\)。这两个条件同时成立意味着 \(\nu+\mu\) 属于 \(\delta(\mathrm{X}_{++})\)；在此情形下，\(\nu+\mu\) 是 \(\mathbf{C}(\nu)\otimes\mathrm{E}(\mu)\) 的最高权。于是，我们对 X(T) 的每个支配权 \(\lambda\) 都构造出了一个最高权为 \(\lambda\) 的 G 的不可约表示，从而在等价意义下得到了 G 的全部不可约表示。由于权为 \(\nu+\mu\) 的向量构成 \(\mathbf{C}(\nu)\otimes\mathrm{E}(\mu)\) 中的一维子空间，证明完成。

推论. 群 G 有一个（有限维）忠实线性表示。

首先注意，X(T) 的每个元素都等于两个支配权之差；更精确地，令 \(\varpi\) 为 \(X_{++}\) 中满足 \(\langle\varpi,K_{\alpha}\rangle>0\) 对每个单根 \(\alpha\) 成立的元素；对于所有 \(\lambda\in X(T)\)，存在正整数 n，使得 \(\langle\lambda+n\varpi,K_{\alpha}\rangle\geq0\) 对每个单根 \(\alpha\) 成立，即（第 1 号，引理 1）\(\lambda+n\varpi\in X_{++}\)。

因此，存在有限族 \((\lambda_{i})_{i\in I}\) 的 \(X_{++}\) 中元素，生成 Z-模 \(\mathrm{X(T)}\)。对于 \(i\in I\)，令 \(\tau_{i}\) 为最高权为 \(\lambda_{i}\) 的 G 的不可约表示（定理 1）；令表示 \(\tau\) 为各 \(\tau_{i}\) 的直和。按构造，权集 \(\mathrm{P}(\tau,\mathrm{T})\) 是 \(\tau\) 的权集，并生成 Z-模 \(\mathrm{X(T)}\)。由第 4 节第 3 号命题 6，现可知同态 \(\tau\) 是单射，故推论得证。

备注。1) 令 \(\mathfrak{n}_{+}\) 为 \(\mathfrak{g}_{\mathbf{C}}\) 的子代数，它是 \(\mathfrak{g}^{\alpha}\) 对所有 \(\alpha > 0\) 的和。令 \(\tau: \mathrm{G} \to \mathrm{GL}(\mathrm{V})\) 为不可约表示，\(\lambda \in \mathrm{X}_{++}\) 为其最高权，\(\tau': \mathfrak{g}_{\mathbf{C}} \to \mathfrak{gl}(\mathrm{V})\) 为由 \(\tau\) 诱导的表示。则 \(\mathrm{V}_{\lambda}(\mathrm{T})\) 是由向量 \(v\) 属于 \(\mathrm{V}\) 且满足 \(\tau'(x)v = 0\) 对所有 \(x \in \mathfrak{n}_{+}\) 成立的向量组成的子空间。

确实，这由 \(\mathfrak{g}_{\mathbf{C}}\)-模 \(\mathbf{C}(\nu) \otimes \mathrm{E}(\mu)\) 的相应断言得到（第 VIII 章，第 6 节，第 2 号，命题 3）。

\begin{enumerate}
\def\labelenumi{\arabic{enumi})}
\setcounter{enumi}{1}
\item
  令 \(\Theta(\mathrm{G})\) 为 \(\mathrm{G}\) 上取值于 \(\mathbf{C}\) 的连续代表函数的代数（《代数》，第 VIII 章）。令 \(\mathrm{G}\) 通过左、右平移作用于 \(\Theta(\mathrm{G})\)。对于每个 \(\lambda \in \mathrm{X}_{++}\)，令 \((\mathrm{V}_{\lambda}, \tau_{\lambda})\) 为 \(\mathrm{G}\) 的最高权为 \(\lambda\) 的不可约表示（定理 1），令 \((\mathrm{V}_{\lambda}^{*}, \check{\tau}_{\lambda})\) 为反步表示（第 III 章，第 3 节，第 11 号）；根据《谱理论》，\(\mathrm{G} \times \mathrm{G}\) 在 \(\Theta(\mathrm{G})\) 上的表示同构于表示 \((\mathrm{V}_{\lambda} \otimes \mathrm{V}_{\lambda}^{*}, \tau_{\lambda} \otimes \check{\tau}_{\lambda})\) 的直和，其中 \(\lambda\) 遍历 \(\mathrm{X}_{++}\)。备注 1 现在推出如下断言：令 \(\lambda \in \mathrm{X}_{++}\)，令 \(\mathrm{E}_{\lambda}\) 为 \(\Theta(\mathrm{G})\) 的向量子空间，它由连续代表函数 \(f\) 组成，这些函数定义在 \(\mathrm{G}\) 上，满足 \(f(gt) = \lambda(t)^{-1} f(g)\) 对所有 \(g \in \mathrm{G}\) 和所有 \(t \in \mathrm{T}\) 成立，并且满足 \(f * x = 0\) 对所有 \(x \in \mathfrak{n}_{-} = \bigoplus_{\alpha < 0} \mathfrak{g}^{\alpha}\) 成立。则 \(\mathrm{E}_{\lambda}\) 在左平移下稳定，且 \(\mathrm{G}\) 通过左平移在 \(\mathrm{E}_{\lambda}\) 上的表示是最高权为 \(\lambda\) 的不可约表示。
\item
  令 \(\tau : \mathrm{G} \to \mathbf{GL}(\mathrm{V})\) 为不可约表示。存在元素 \(\nu\) 属于 \(\mathrm{X}(\mathrm{C}(\mathrm{G}))\)，使得 \(\tau(s)v = \nu(s)v\) 对所有 \(s \in \mathrm{C}(\mathrm{G}), v \in \mathrm{V}\) 成立：事实上，\(\tau(\mathrm{C}(\mathrm{G}))\) 包含于 \(\tau(\mathrm{G})\) 的交换子中，而后者等于 \(\mathbf{C}^*.1_{\mathrm{V}}\)（《代数》，第 VIII 章，第 3 节，第 2 号，定理 1）。对于每个权 \(\lambda\) 属于 \(\tau\)，\(\lambda\) 在 \(\mathrm{C}(\mathrm{G})\) 上的限制等于 \(\nu\)。
\item
  第 VIII 章第 7 节第 2 号至第 5 号的定义和断言可以毫无困难地推广到当前情形；细节留给读者。
\end{enumerate}

命题 1. 设 \(\tau : G \to \mathbf{GL}(V)\) 是 \(G\) 的最高权为 \(\lambda \in X_{++}\) 的不可约表示。令 \(m\) 为整数 \(\sum_{\alpha \in R_+} \langle \lambda, K_\alpha \rangle\)，令 \(w_0\) 为满足 \(w_0(R_+) = R_-\) 的 Weyl 群元素（第 VI 章，第 6 节，命题 17 的推论 3）。有三种可能的情形：

\begin{enumerate}
\def\labelenumi{\alph{enumi})}
\item
  \(w_{0}(\lambda) = -\lambda\) 且 m 为偶数。则 V 上存在 G-不变的分离对称双线性型；表示 \(\tau\) 为实型（附录 II）。
\item
  \(w_{0}(\lambda) \neq -\lambda\)。V 上每个 G-不变双线性型都为零；表示 \(\tau\) 为复型（同上引文）。
\item
  \(w_{0}(\lambda) = -\lambda\) 且 m 为奇数。则 V 上存在 G-不变的分离交错双线性型；表示 \(\tau\) 为四元数型（同上引文）。
\end{enumerate}

若 \(\tau\) 在 \(\mathrm{C}(\mathrm{G})_0\) 上的限制不是平凡的，则属于情形 \(b\)。

V 上的双线性型 B 在 G 下不变，当且仅当它在 \(g_{C}\) 下不变（第 III 章，第 6 节，第 5 号，推论 3）。因此，若 G 为半单的，则命题由第 VIII 章第 7 节第 5 号命题 12 以及附录 II 的命题 3 得出。

在一般情形下，置 \(\mathrm{C}(\mathrm{G})_{0} = \mathrm{S}\)，并将 \(\mathrm{X}(\mathrm{T}/\mathrm{S})\) 与 \(\mathrm{X}(\mathrm{T})\) 的一个子群识别（该子群在 W 下稳定）。若 \(\tau(\mathrm{S}) = \{1_{\mathrm{V}}\}\)，则 \(\tau\) 通过取商诱导出 G/S 的表示 \(\tau' : G/S \to GL(V)\)，其最高权为 \(\lambda\)；在此情形下，将前面的结果应用于 G/S 即得命题。

假设 \(\tau (\mathrm{S})\neq \{1_{\mathrm{V}}\}\)。存在非零元素 \(\nu\) 属于 \(\mathrm{X(S)}\)，使得 \(\tau (s) = \nu (s)_{\mathrm{V}}\) 对所有 \(s\in \mathrm{S}\) 成立（备注 3）。于是 \(\nu\) 是 \(\lambda\) 在限制同态 \(\mathrm{X}(\mathrm{T})\to\mathrm{X}(\mathrm{S})\) 下的像；由于 W 在 \(\mathrm{X}(\mathrm{S})\) 上平凡作用，等式 \(w_{0}(\lambda)=-\lambda\) 蕴含 \(\nu=-\nu\)，这是不可能的；所以 \(w_{0}(\lambda)\neq-\lambda\)。另一方面，若 B 是 V 上 G-不变的双线性型，则对于 V 中所有 x,y 以及 S 中 s，有

\[
\mathrm{B} (\nu (s) x, \nu (s) y) = \mathrm{B} (x, y) = \nu (s) ^ {2} \mathrm{B} (x, y)
\]

这蕴含 B = 0，故命题得证。

令 \(\mathfrak{S}_{\mathbf{R}}(\mathrm{G})\) 为 G 在有限维实向量空间上的不可约连续表示的等价类集合。命题 1 及附录 II 的结果给出双射 \(\Phi:X_{++}/\Sigma\to\mathfrak{S}_{\mathbf{R}}(\mathrm{G})\)，其中 \(\Sigma\) 表示子群 \(\{1,-w_{0}\}\) \(\operatorname{Aut}(\mathrm{X}(\mathrm{T}))\) 的子群。更精确地，令 \(\lambda\in X_{++}\)，令 \(E_{\lambda}\) 为最高权为 \(\lambda\) 的 G 的一个表示；则

\[
\begin{array}{c} \Phi (\{\lambda , - w _ {0} (\lambda) \}) = \mathrm{E} _ {\lambda [ \mathbf {R} ]} \text {if} \lambda \neq - w _ {0} (\lambda) \text {or if} \sum_ {\alpha \in \mathrm{R} _ {+}} \langle \lambda , K _ {\alpha} \rangle \notin 2 \mathbf {Z} \\ \Phi (\{\lambda , - w _ {0} (\lambda) \}) = \mathrm{E} _ {\lambda} ^ {\prime} \text {if} \lambda = - w _ {0} (\lambda) \text {and} \sum_ {\alpha \in \mathrm{R} _ {+}} \langle \lambda , K _ {\alpha} \rangle \in 2 \mathbf {Z} \end{array}
\]

其中 \(E_{\lambda}^{\prime}\) 是 \(E_{\lambda}\) 上在 G 下不变的 R-结构。

